% Apendice A (revisao pos-EE). Gerado por 02_audit/build12_appendix.py; tabulars verbatim.
\begin{table}[htbp]\centering
\caption{Birth-cohort applicants and the duration sample}
\label{tab:P2rec}
\footnotesize
\begin{tabular}{lrrrr}
\toprule
Cohort & Total & Duration & Excl.$>$47m & \% excl. \\
\midrule
2010 & 128,352 & 106,300 & 20,322 & 15.8 \\
2011 & 146,693 & 126,766 & 18,556 & 12.6 \\
2012 & 155,292 & 139,440 & 14,638 & 9.4 \\
2013 & 153,033 & 140,811 & 11,178 & 7.3 \\
2014 & 151,772 & 146,665 & 4,144 & 2.7 \\
2015 & 146,293 & 146,232 & 0 & 0.0 \\
2016 & 128,143 & 128,143 & 0 & 0.0 \\
\bottomrule
\end{tabular}
\tabnote{\emph{Source:} SME childcare request record, 2010--2018 applications, with referral and enrollment dates through November 2019. \emph{Sample:} birth cohorts 2010--2016. \emph{Variables:} Total: children with a valid first application; Duration: those with non-negative eligible time; Excl.$>$47m: children whose age at first application exceeds 47 whole months. \emph{Exercise:} counts. Adding the excluded children as immediate aging out widens the referral gain (Table~\ref{tab:P2rob}). Total minus Duration (22,052 in 2010) exceeds Excl.$>$47m (20,322) because Table~\ref{tab:agingout} dates the limit by day (application after the day the child turns 47 months), whereas this column uses age at application in whole months; the day-based count is the one used in Table~\ref{tab:agingout}.}
\end{table}
